\documentclass[UTF8,12pt]{ctexart}
\usepackage{amsmath}
\usepackage{geometry}
\geometry{a4paper, margin=2.5cm}

\title{中文文档：自动切换 XeLaTeX}
\author{测试作者}
\date{\today}

\begin{document}
\maketitle

\section{引言}
这个文件没有魔法注释，也没有手选引擎。TeXMini 应该根据 \verb|ctexart| 自动使用 XeLaTeX，状态栏显示 \texttt{latexmk · xelatex}。

如果状态栏显示的是 pdflatex，或者编译出来是乱码 / 报错 \verb|! Package ctex Error|，说明启发式失效了。

\section{中英混排与标点}
中文、English、数字 123、全角标点，。！？；：以及半角标点 ,.!?;: 混排。
行内公式 $\alpha + \beta = \gamma$，行间公式：
\begin{equation}
  \int_0^1 x^2 \,\mathrm{d}x = \frac{1}{3}
\end{equation}

\section{字数统计}
状态栏的"词"数对中文是按字统计还是按空格分词？这一段有明确的字数：一二三四五六七八九十。
观察状态栏数字是否随输入变化，以及是否合理（不应把整段中文算作 1 个词）。

\subsection{大纲里的中文标题}
大纲面板应正确显示中文章节名并且点击可跳转。

\section{SyncTeX 与中文}
在这一行双击，PDF 应定位到"SyncTeX 与中文"这一节。反过来在 PDF 里双击这一段，编辑器应回到这里而不是偏到别的行。

\end{document}
